\documentclass[a4paper,12pt]{article}
\usepackage[T2A]{fontenc}
\usepackage[utf8]{inputenc}
\usepackage[russian]{babel}
\usepackage{amsmath,amssymb,mathtools}
\usepackage{helvet}
\renewcommand{\sfdefault}{cmss}
\renewcommand{\familydefault}{\sfdefault}
\usepackage{icomma}
\usepackage[left=23mm,right=23mm,top=22mm,bottom=23mm,headheight=26pt]{geometry}
\usepackage{microtype}
\usepackage{tikz}
\usetikzlibrary{calc,arrows.meta,positioning,shapes.geometric}
\usepackage{tabularx,booktabs,array}
\usepackage{enumitem}
\usepackage[hidelinks]{hyperref}
\usepackage{parskip}
\usepackage{fancyhdr}
\usepackage{setspace}
\usepackage{xcolor}
\usepackage{titlesec}
\usepackage{needspace}
\usepackage{float}

\definecolor{accent}{HTML}{008080}
\definecolor{darkaccent}{HTML}{004D4D}
\definecolor{textgray}{HTML}{303638}
\definecolor{softgray}{HTML}{6B7576}
\color{textgray}
\setstretch{1.39}
\setlength{\parindent}{0pt}
\setlength{\parskip}{3pt}
\setlist[itemize]{leftmargin=6mm,itemsep=1pt,topsep=4pt,parsep=0pt}
\setlist[enumerate]{leftmargin=8mm,itemsep=9pt,topsep=8pt,parsep=1pt}
\titleformat{\section}{\large\color{darkaccent}\bfseries}{\thesection.}{.65em}{}
\titlespacing*{\section}{0pt}{16pt}{8pt}
\titleformat{\subsection}{\normalsize\color{accent}\bfseries}{}{0pt}{}
\titlespacing*{\subsection}{0pt}{12pt}{5pt}
\setlength{\headsep}{10pt}
\pagestyle{fancy}
\fancyhf{}
\fancyhead[L]{\small\color{softgray}Теория вероятностей}
\fancyhead[R]{\small\color{softgray}Марк Гукин}
\renewcommand{\headrulewidth}{0.45pt}
\renewcommand{\headrule}{\hbox to\headwidth{\color{accent}\leaders\hrule height \headrulewidth\hfill}}
\fancyfoot[C]{\small\color{softgray}\thepage}
\renewcommand{\footrulewidth}{0pt}
\newcommand{\keyidea}[1]{\textcolor{darkaccent}{\bfseries #1}}
\newcommand{\ex}[1]{\par\smallskip\textcolor{accent}{\bfseries #1}\par\smallskip}
\newcommand{\ordpair}[2]{(#1,#2)}
\setcounter{secnumdepth}{1}
\setcounter{tocdepth}{1}
\year=2026\month=11\day=19

\begin{document}
\begin{titlepage}
\thispagestyle{empty}
\centering
\vspace*{26mm}
{\normalsize\color{darkaccent}\MakeUppercase{Математика}\par}
\vspace{18mm}
{\LARGE Чек-лист\par}
\vspace{16mm}
{\LARGE\color{accent} Теория вероятностей\\[3mm]
\large Кубик, монета, комбинаторика\\ и противоположные события\par}
\vspace{19mm}
{\normalsize Подготовка к ЕГЭ по профильной математике\par}
\vfill
{\large Марк Гукин\par}
\vspace{12mm}
{\normalsize \today\par}
\vspace{20mm}
{\small Лёвин Артём Александрович\\[2mm]
эксклюзивно для Марка Гукина\par}
\vspace{21mm}
\begin{tikzpicture}[x=1cm,y=1cm]
 \draw[accent,line width=0.85pt] (0,0) .. controls (1.0,.38) and (2.2,-.38) .. (3.4,0)
 .. controls (4.6,.35) and (5.6,-.35) .. (6.7,0)
 .. controls (7.85,.33) and (8.9,-.32) .. (10,0);
\end{tikzpicture}
\vspace{3mm}
\end{titlepage}

\section{Вероятность и равновозможные исходы}
Если все элементарные исходы равновозможны, вероятность события $A$ равна
\[
\boxed{\displaystyle P(A)=\frac{N(A)}{N}},
\]
где $N$~--- число всех исходов, $N(A)$~--- число благоприятных исходов.

\keyidea{Проверяйте слова в условии:}
\begin{itemize}
\item \emph{меньше четырёх} означает $1,2,3$; четвёрка исключена;
\item \emph{не более четырёх} означает $1,2,3,4$;
\item \emph{ровно два} означает исключительно два;
\item \emph{хотя бы два} (\emph{не менее двух}) означает два или больше.
\end{itemize}
\ex{Пример 1. Один правильный кубик}
Вероятность выпадения числа меньше четырёх:
\[
P=\frac{3}{6}=\frac12.
\]
\textcolor{softgray}{Слова \emph{больше}, \emph{меньше}, \emph{не более}, \emph{не менее} определяют, какие граничные значения нужно учитывать.}

\section{Метод позиций: последовательные броски}
\keyidea{Учитывайте порядок испытаний.} Два броска кубика задают упорядоченную пару $(a,b)$; например, $(2,5)$ и $(5,2)$~--- разные исходы. Для каждого броска имеется по шесть вариантов:
\[
N=\underbrace{6\cdot6}_{\text{два броска}}=36.
\]
При трёх бросках монеты, обозначая орла буквой $\mathrm{O}$, а решку~--- $\mathrm{P}$, получаем $2^3=8$ исходов.

\ex{Пример 2. Два одинаковых числа на двух кубиках}
Благоприятные пары: $(1,1),(2,2),\ldots,(6,6)$. Их шесть, поэтому
\[
P=\frac{6}{36}=\frac16.
\]
\ex{Пример 3. Ровно два орла за три броска}
Подходят $\mathrm{OOP}$, $\mathrm{OPO}$, $\mathrm{POO}$~--- три исхода из восьми:
\[
P=\frac38.
\]

\section{Когда нужна таблица исходов}
При двух бросках кубика удобно построить таблицу сумм; строка задаёт первый бросок, столбец~--- второй.
\begin{center}
\small\renewcommand{\arraystretch}{1.19}
\begin{tabular}{c@{\hspace{6pt}}*{6}{>{\centering\arraybackslash}p{1.05cm}}}
\toprule
$+ $ & $1$ & $2$ & $3$ & $4$ & $5$ & $6$ \\
\midrule
$1$ & $2$ & $3$ & $4$ & $5$ & $6$ & $7$ \\
$2$ & $3$ & $4$ & $5$ & $6$ & $7$ & $8$ \\
$3$ & $4$ & $5$ & $6$ & $7$ & $8$ & $9$ \\
$4$ & $5$ & $6$ & $7$ & $8$ & $9$ & $10$ \\
$5$ & $6$ & $7$ & $8$ & $9$ & $10$ & $11$ \\
$6$ & $7$ & $8$ & $9$ & $10$ & $11$ & $12$ \\
\bottomrule
\end{tabular}
\end{center}
Чётная сумма получается, когда оба числа имеют одинаковую чётность. Таких клеток $18$ из $36$, следовательно, $P=\frac{18}{36}=\frac12$.

\section{Умножение и сложение вероятностей}
Для \emph{независимых} событий $A$ и $B$:
\[
P(A\cap B)=P(A)P(B).
\]
Если необходимо, чтобы состоялись \emph{оба} события, в условиях независимости их вероятности умножают (союз «и»).

Если исходы $A$ и $B$ \emph{несовместны}, то
\[
P(A\cup B)=P(A)+P(B).
\]
Значит, для выбора между \emph{непересекающимися} случаями (союз «или») вероятности складывают. Если случаи пересекаются, просто складывать нельзя.

\ex{Пример 4. Ровно одна бракованная ручка}
Из двух независимо выбранных ручек каждая бракована с вероятностью $0{,}1$. Благоприятны два несовместных случая: первая бракована, вторая исправна; или наоборот.
\[
P=0{,}1\cdot0{,}9+0{,}9\cdot0{,}1=0{,}18.
\]

\section{Противоположное событие: «хотя бы один»}
Для события $A$ и противоположного события $\overline A$:
\[
\boxed{P(A)=1-P(\overline A).}
\]
\keyidea{«Хотя бы один успех»} означает, что противоположное событие~--- \emph{ни одного успеха}.

\ex{Пример 5. Монету бросают трижды}
Вероятность хотя бы одного орла удобнее вычислить через три решки:
\[
P(\text{хотя бы один орёл})
=1-P(\mathrm{PPP})=1-\left(\frac12\right)^3=\frac78.
\]
Если независимое испытание повторяют $n$ раз, а успех в каждом испытании имеет вероятность $p$, то
\[
\boxed{P(\text{хотя бы один успех})=1-(1-p)^n.}
\]
В частности, для пяти бросков кубика вероятность \emph{ни разу} не получить тройку равна $(5/6)^5$, а \emph{получить её хотя бы раз}~--- $1-(5/6)^5$.

\vspace{4mm}
\subsection*{Самопроверка: четыре частые ошибки}
\begin{itemize}
\item Путают \emph{ровно один} и \emph{хотя бы один}: второй случай включает и несколько успехов.
\item Складывают вероятности пересекающихся событий как несовместные.
\item Умножают вероятности событий, не проверив независимость.
\item Смешивают выбор \emph{группы} и порядок выступлений участников.
\end{itemize}


\clearpage
\begin{center}
{\Large\color{darkaccent}\bfseries Алгоритм: «хотя бы один успех»}\par
\vspace{7mm}
\end{center}
\tikzset{
 fcStart/.style={ellipse,draw=accent,line width=.75pt,minimum width=5.2cm,minimum height=1cm,align=center,font=\small},
 fcAction/.style={rectangle,rounded corners=2pt,draw=accent,line width=.75pt,text width=8.2cm,minimum height=1.16cm,align=center,font=\small,inner sep=6pt},
 fcDecision/.style={diamond,aspect=2.45,draw=accent,line width=.75pt,text width=4.6cm,align=center,font=\small,inner sep=2pt},
 fcArrow/.style={-{Stealth[length=2.6mm]},draw=darkaccent,line width=.8pt}
}
\begin{center}
\begin{tikzpicture}[node distance=12mm and 12mm]
 \node[fcStart] (st) {Начало: требуется вероятность\\ хотя бы одного успеха};
 \node[fcAction,below=of st] (neg) {Перейдите к противоположному событию: \\ «успех не произошёл ни разу»};
 \node[fcDecision,below=17mm of neg] (ind) {Все испытания независимы?};
 \node[fcAction,text width=5.25cm,below left=20mm and 3mm of ind] (no) {Найдите $P(\overline A)$\\перебором или другим\\обоснованным способом};
 \node[fcAction,text width=5.25cm,below right=20mm and 3mm of ind] (yes) {Найдите $P(\overline A)$\\как произведение вероятностей\\неудачи во всех испытаниях};
 \node[fcAction,below=41mm of ind] (sub) {Вычтите вероятность противоположного события:\\ $P(A)=1-P(\overline A)$};
 \node[fcStart,below=of sub] (end) {Ответ};
 \draw[fcArrow] (st) -- (neg);
 \draw[fcArrow] (neg) -- (ind);
 \draw[fcArrow] (ind.south west) -| node[pos=.28,above,sloped,fill=white,inner sep=1pt]{нет} (no.north);
 \draw[fcArrow] (ind.south east) -| node[pos=.28,above,sloped,fill=white,inner sep=1pt]{да} (yes.north);
 \draw[fcArrow] (no.south) |- ([yshift=5mm]sub.west) -- (sub.west);
 \draw[fcArrow] (yes.south) |- ([yshift=5mm]sub.east) -- (sub.east);
 \draw[fcArrow] (sub) -- (end);
\end{tikzpicture}
\end{center}
\clearpage

\section{Выбор людей и порядок выступлений}
При выборе нескольких людей \emph{без учёта порядка} перестановки выбранных участников не создают новых групп. Из $n$ различных людей можно выбрать $k$ человек числом способов
\[
\binom nk=\frac{n!}{k!(n-k)!},\qquad 0\le k\le n.
\]
\ex{Пример 6. Выбор трёх туристов из двадцати четырёх}
Нужно, чтобы выбранный турист $A$ обязательно оказался в группе. Его место занято; оставшихся двоих выбираем из $23$:
\[
P=\frac{\binom{23}{2}}{\binom{24}{3}}=\frac{3}{24}=\frac18.
\]

При случайном определении \emph{порядка} выступлений фиксированная позиция равновероятна для каждого участника. Если спортсменов всего $20$, из которых $7$~--- из России, то вероятность увидеть спортсмена из России на любом конкретном месте (первом, пятом, последнем) равна
\[
P=\frac{7\cdot19!}{20!}=\frac7{20}.
\]
\textcolor{softgray}{Число способов получить группу и число способов выстроить очередь различаются: до подсчёта выясните, существенен ли порядок.}

\section{Промежутки и вложенные события}
Если $a<b$, то события $\{X<a\}$ и $\{a\le X<b\}$ несовместны и вместе образуют событие $\{X<b\}$. Поэтому
\[
\boxed{P(a\le X<b)=P(X<b)-P(X<a).}
\]
\begin{center}
\begin{tikzpicture}[x=0.112cm,y=1cm]
 \draw[->,line width=.65pt] (0,0) -- (96,0);
 \draw[accent,line width=2.4pt] (31,0) -- (72,0);
 \draw[accent,fill=white,line width=1pt] (31,0) circle (2.3pt);
 \draw[accent,fill=white,line width=1pt] (72,0) circle (2.3pt);
 \draw[darkaccent,fill=darkaccent] (31,0) circle (1.45pt);
 \node[below=6pt,font=\small] at (31,0) {$a$};
 \node[below=6pt,font=\small] at (72,0) {$b$};
 \node[above=8pt,font=\small,accent] at (51.5,0) {$a\le X<b$};
\end{tikzpicture}
\end{center}
\ex{Пример 7. Масса упаковки}
Пусть $P(X<810)=0{,}96$ и $P(X<790)=0{,}07$. Тогда
\[
P(790\le X<810)=0{,}96-0{,}07=0{,}89.
\]

\clearpage
\section*{Тренировочный блок}
\textcolor{softgray}{Путь А · 10 задач · от базовых к комплексным. В каждом задании запишите число всех исходов или обоснуйте применённую формулу. Ответы приведены отдельно.}
\subsection*{Базовый уровень}
\begin{enumerate}[label=\textcolor{darkaccent}{\bfseries \arabic*.},ref=\arabic*]
\item Правильный игральный кубик бросают один раз. Найдите вероятность того, что выпавшее число \emph{больше двух, но не более пяти}.
\item Правильный игральный кубик бросают дважды. Найдите вероятность того, что сумма выпавших чисел равна $7$.
\item Симметричную монету бросают трижды. Найдите вероятность того, что орёл выпадет \emph{ровно два раза}.
\end{enumerate}
\subsection*{Уверенное владение}
\begin{enumerate}[label=\textcolor{darkaccent}{\bfseries \arabic*.},start=4,ref=\arabic*]
\item Правильный игральный кубик бросают дважды. Найдите вероятность того, что сумма выпавших чисел равна $8$. Для проверки можете составить таблицу исходов.
\item Симметричную монету бросают четыре раза. Найдите вероятность того, что орёл выпадет \emph{не менее трёх раз}.
\item Из группы $15$ туристов случайно выбирают троих. Найдите вероятность того, что в эту тройку войдёт конкретный турист $A$. Порядок выбора роли не играет.
\item На соревновании выступают $20$ спортсменов, из них $6$ представляют одну страну. Порядок выступлений определяют случайно. Найдите вероятность того, что на \emph{двенадцатой} позиции окажется спортсмен этой страны.
\end{enumerate}
\subsection*{Комплексное применение}
\begin{enumerate}[label=\textcolor{darkaccent}{\bfseries \arabic*.},start=8,ref=\arabic*]
\item Для массы упаковки $X$ известно: $P(X<850)=0{,}91$, а $P(X<800)=0{,}14$. Найдите вероятность события $800\le X<850$.
\item Две ручки выбирают независимо. Вероятность брака каждой ручки равна $0{,}1$. Найдите вероятность того, что \emph{ровно одна} из двух ручек окажется бракованной.
\item Правильный игральный кубик бросают пять раз. Найдите вероятность того, что тройка выпадет \emph{хотя бы один раз, но не во всех пяти бросках}. Ответ представьте обыкновенной дробью.
\end{enumerate}

\clearpage
\section*{Ответы}
\vspace{5mm}
\renewcommand{\arraystretch}{1.7}
\begin{center}
\begin{tabularx}{.91\linewidth}{>{\centering\arraybackslash}p{.8cm}X>{\centering\arraybackslash}p{.8cm}X}
\toprule
\textbf{№} & \textbf{Ответ} & \textbf{№} & \textbf{Ответ} \\
\midrule
1 & $\displaystyle \frac12$ & 6 & $\displaystyle \frac15$ \\[8pt]
2 & $\displaystyle \frac16$ & 7 & $\displaystyle \frac3{10}$ \\[8pt]
3 & $\displaystyle \frac38$ & 8 & $0{,}77$ \\[8pt]
4 & $\displaystyle \frac5{36}$ & 9 & $0{,}18$ \\[8pt]
5 & $\displaystyle \frac5{16}$ & 10 & $\displaystyle \frac{775}{1296}$ \\[8pt]
\bottomrule
\end{tabularx}
\end{center}
\vspace{12mm}
\section*{Итоговая памятка}
\renewcommand{\arraystretch}{1.4}
\small
\begin{tabularx}{\linewidth}{>{\raggedright\arraybackslash}p{.43\linewidth}X}
\toprule
\textbf{Формулировка} & \textbf{Полезный приём} \\
\midrule
Ровно $k$ успехов & Учитывайте только исходы с $k$ успехами. \\
Хотя бы один успех & Вычтите из единицы вероятность всех неудач. \\
И (независимые события) & Перемножьте вероятности. \\
Или (несовместные случаи) & Сложите вероятности. \\
Выбор без учёта порядка & Используйте сочетания. \\
Фиксированное место в случайной очереди & Посчитайте, сколько подходящих участников из всех. \\
Вероятность промежутка & Вычтите вложенные накопленные вероятности. \\
\bottomrule
\end{tabularx}
\normalsize
\end{document}
